% Ukrainian translation of Correll et al., Introduction to Autonomous Robots
% Source: github.com/Correll-Lab/Introduction-to-Autonomous-Robots
% File: chapters/vision.tex
% CC BY-NC-SA license per repo.
% Translation by auxiliary local run (Anthropic), 2026-06-01.
\chapter{Зір (Vision)}\label{chap:vision}

Зір — одна з найінформативніших сенсорних систем, доступних і людям, і роботам. Однак ефективна та точна обробка цієї багатої інформації досі залишається ключовим викликом галузі. Цілі цього розділу:
\begin{itemize}
\item ввести поняття зображення як двовимірного сигналу;
\item дати інтуїцію щодо багатства інформації, прихованої в низькорівневих даних;
\item ввести базові алгоритми обробки зображень — згорткові та порогові.
\end{itemize}

\section{Зображення як двовимірні сигнали}

Зображення захоплюються камерами, які містять матриці приладів зі зарядовим зв'язком (CCD) або подібних напівпровідників (наприклад, КМОП, CMOS), що перетворюють фотони на електричні сигнали. Ці матриці зчитуються попіксельно і перетворюються на цифрові значення, наприклад масив $640 \times 480$ трибайтових кортежів RGB. Дивлячись на дані у матриці, чітко видно білу плитку серед чорних у нижньому правому куті шахівниці. Вищі значення відповідають яскравішим кольорам (білий), нижчі — темнішим. Хоча плитки мають однаковий колір, фактичні значення помітно відрізняються. Має сенс думати про ці значення як про 1D-сигнал: ``похідна'' уздовж горизонтальних рядків укаже на області різких змін, а ``частотна'' гістограма зображення вкаже, як швидко змінюються значення. Області з плавними градієнтами (чорно-білі плитки) матимуть низькі частоти; області з сильними градієнтами — високочастотну інформацію.

\begin{figure}
    \centering
    \includegraphics[width=\textwidth]{figs/iss_closeupmatrix}
    \caption{Шахівниця всередині МКС з астронавтом Грегорі Хамітофф. У вставці — приклад фактичних даних сенсора. Чітко видно контури білої плитки.}
    \label{fig:iss_closeup}
\end{figure}

Ця мова відкриває серії концепцій сигналів: \emph{фільтри нижніх частот} (придушують високочастотну інформацію), \emph{фільтри верхніх частот} (придушують низькочастотну), \emph{смугові фільтри} (пропускають лише діапазон), аналіз частотного спектру, ``згортання'' зображення з іншою двовимірною функцією.

\section{Від сигналів до інформації}

Багато явищ з різним або навіть протилежним змістом виглядають подібно на низькому рівні. Наприклад, різкі зміни кольору не обов'язково означають зміну кольору поверхні. Подібні патерни виникають також від розривів глибини, дзеркальних відблисків, змін освітлення або змін орієнтації поверхні. Це — те, що робить комп'ютерний зір складною задачею.

\begin{figure}[!htb]
    \centering
    \includegraphics[width=\textwidth]{figs/iss_edges}
    \caption{Всередині МКС (зліва); області з різкими змінами пікселів (справа). Причини: зміна властивостей поверхні (1), розриви глибини (2), дзеркальні відблиски (3), зміни освітлення/тіні (4), зміни орієнтації поверхні (5).
    \label{fig:iss_edges}}
\end{figure}

Цей приклад ілюструє: сигнали й дані самі по собі недостатні для розуміння — потрібен \emph{контекст}. Контекст включає не лише сусідні сигнали, а й високорівневі концептуальні знання: те, що джерела світла створюють тіні й дзеркальні відблиски, що передні об'єкти здаються більшими, тощо. Важливість концептуальних знань ілюструє рис.~\ref{fig:craters}: обидва зображення показують ідентичний ландшафт, але один здається покритим кратерами, інший — пагорбами. Знаючи, що Сонце освітлює сцену з лівого/правого боку, парадокс прояснюється — змінне освітлення робить кратери схожими на горби.

\begin{figure}[!htb]
    \centering
    \includegraphics[width=\textwidth]{figs/craters}
    \caption{Знімок місця посадки Apollo 15 у різний час доби. Ландшафт однаковий, але здається покритим кратерами (зліва) або пагорбами (справа). Знаючи напрямок Сонця, ефект прояснюється. Image credit: NASA/GSFC/Arizona State University.
    \label{fig:craters}}
\end{figure}

Більше того, концептуальних знань часто достатньо, щоб компенсувати брак низькорівневих ознак. Приклад на рис.~\ref{fig:dalmatian} — далматинця можна впізнати, попри відсутність більшості його контурів.

Ці приклади ілюструють переваги й недоліки сигнально-обробного підходу. Алгоритм детектує цікаві сигнали навіть там, де ми їх не бачимо чи не очікуємо (через концептуальне зміщення). Розуміння зображення вимагає не лише низькорівневої обробки, а й розумного поєднання просторових відношень з концептуальними знаннями. Це робиться через згорткові нейромережі (\cref{chap:ann}).

\begin{figure}[!]
    \centering
    \includegraphics[width=\textwidth]{figs/dalmatian}
    \caption{Зображення далматинця чітко розпізнається попри відсутність низькорівневих ознак (краї лише для вух, підборіддя, частини лап).
    \label{fig:dalmatian}}
\end{figure}

\section{Базові операції з зображеннями}

Базові операції — це фільтри в частотній або просторовій (інтенсивність/колір) області. Зазвичай фільтр діє в інтенсивностній області, але знання його впливу на частоту допомагає розуміти його функцію. Фільтр для виділення країв повинен пригнічувати низькі частоти (області, де значення не змінюються) і підсилювати високі (де змінюються швидко).

\subsection{Порогові операції}
Щоб знайти об'єкти певного кольору або амплітуди краю, \emph{порогування} зображення дає бінарне зображення з ``true/false'' регіонами. Використовуються оператори $>, <, \leq, \geq$. Існують і адаптивні версії, що локально оновлюють пороги (для змін освітлення).
Хоча порогування просте, знаходження правильних значень — складна задача: значення пікселів сильно змінюються з освітленням, тож немає такого як ``червоний'' або ``зелений'' під різними умовами.

\subsection{Згорткові фільтри}

Фільтр реалізується оператором \emph{згортки}\index{згортка} $\star$, що згортає функцію $f()$ з функцією $g()$:
\begin{equation}
f(x)\star g(x) = \int_{-\infty}^{\infty} f(\tau) g(x-\tau)\, d\tau,
\end{equation}
де $g()$ — \emph{фільтр}. Згортка ``ковзає'' $g()$ вздовж $f()$, множачи їх. Оскільки зображення дискретні, згортка теж дискретна:
\begin{equation}
f[x]\star g[x] = \sum_{i=-\infty}^{\infty} f[i] g[x-i].
\end{equation}
Для 2D-зображень — двовимірна:
\begin{equation}\label{eq:2dconv1}
f[x,y]\star g[x,y] = \sum_{i=-\infty}^{\infty}\sum_{j=-\infty}^{\infty} f[i,j] g[x-i, y-j].
\end{equation}
Хоча визначення від $-\infty$ до $\infty$, зображення і фільтри скінченні. Згортка комутативна:
\begin{equation}\label{eq:2dconv2}
f[x,y]\star g[x,y] = \sum_{i,j} f[x-i, y-j] g[i,j].
\end{equation}

\subsubsection{Гаусове згладжування}

Один із найважливіших фільтрів — \emph{гаусів фільтр}\index{гаусів фільтр}, що має форму гаусівського дзвону:
\begin{equation}
g(x,y) = \frac{1}{10}
\left(
\begin{array}{ccc}
1 & 1 & 1\\
1 & 2 & 1\\
1 & 1 & 1
\end{array}
\right)
\end{equation}
Кожен піксель стає середнім сусідів, з власним значенням, зваженим вдвічі більше. Гаусів фільтр діє як \emph{фільтр нижніх частот}\index{фільтр нижніх частот} — пригнічує високі частоти. Шум у зображенні згладжується.

\begin{figure}[!]
    \centering
    \includegraphics[width=\textwidth]{figs/filters}
    \caption{Зашумлене зображення до (вгорі ліворуч) і після гаусового фільтра (вгорі праворуч). Відповідні зображення країв — нижче.
    \label{fig:filters}}
\end{figure}

\subsubsection{Детектування країв}\label{sec:sobel}

Детектування країв — згорткою з ядром Собеля\index{фільтр Собеля}:
\begin{equation}
s_x = \left(\begin{array}{ccc} -1 & 0 & 1\\ -2 & 0 & 2\\ -1 & 0 & 1 \end{array}\right), \qquad
s_y = \left(\begin{array}{ccc} 1 & 2 & 1\\ 0 & 0 & 0\\ -1 & -2 & -1 \end{array}\right).
\end{equation}
$s_x$ детектує вертикальні краї, $s_y$ — горизонтальні. \emph{Детектор країв Кенні} (Canny)\index{детектор Кенні} запускає принаймні два таких фільтри.

\subsubsection{Різниця гаусіан (DoG)}\label{sec:vision:dog}

\emph{Difference of Gaussians} (DoG)\index{Difference of Gaussians}: віднімання двох зображень, кожне відфільтроване гаусіаном різної ширини. Обидва пригнічують високочастотну інформацію; їхня різниця — \emph{смуговий фільтр}, з якого вилучено і низькі, і високі частоти. Одне ядро зазвичай у 4--5 разів ширше за інше.

DoG також апроксимує \emph{лапласіан гаусіана}\index{лапласіан гаусіана} (Laplacian of Gaussian, LoG) — суму других похідних гаусівського ядра. Одне ядро приблизно в 1.6 разів ширше за інше. DoG/LoG виокремлюють високочастотну інформацію (краї), пригнічуючи високочастотний шум.

\subsection{Морфологічні операції}

Інший клас фільтрів — \emph{морфологічні оператори} з ядром, що описує структуру операції, і правилом зміни пікселя за околом. Важливі: \emph{ерозія}\index{ерозія} (мінімум у околі) і \emph{дилатація}\index{дилатація} (максимум у околі). Корисні для заповнення дірок у лінії або видалення шуму. Дилатація + ерозія = ``Закриття'' (Closing); ерозія + дилатація = ``Відкриття'' (Opening). Віднімання еродованого і дилатованого зображень — детектор країв.

\begin{figure}
    \centering
    \includegraphics[width=\textwidth]{figs/morphology}
    \caption{Приклади морфологічних операторів — ерозія, дилатація та комбінації (з документації OpenCV, BSD).
    \label{fig:morphology}}
\end{figure}

\section{Виокремлення структури із зору}\label{sec:vision:structure}

Зір дає і \emph{семантичну} (якості — що в сцені), і \emph{метричну} (кількості — розміри, відстані) інформацію. Семантика нині сильно покладається на машинне навчання. Метрика — на геометричних відношеннях.

Рис.~\ref{fig:stereovision} показує відношення між двома кадрами, що спостерігають одну точку. Не розрізняємо тут просторово- і часово-корельовані випадки — це дві різні задачі: \emph{стереозір}\index{стереозір} (дві жорстко прикріплені камери, просторово корельовані); \emph{структура з руху}\index{структура з руху} (одна камера переміщується, пара зображень корельована через матрицю перетворення). У стереозорі це перетворення — \emph{екстринсики сенсора} (6-DoF, калібровані). У структурі з руху — рух камери, що оцінюється локалізацією.

\begin{figure}
\centering
    \def\svgwidth{1.0\textwidth}
    \import{./figs/}{stereovision.pdf_tex}
\caption{Схема кореляції ознак між зображеннями для виокремлення тривимірної інформації з двовимірних видів.}
\label{fig:stereovision}
\end{figure}

3D-точка проектується в кадр камери:
\begin{equation}
\begin{pmatrix} u\\ v\\ 1 \end{pmatrix}
= K T \begin{pmatrix} x\\ y\\ z\\ 1 \end{pmatrix},
\label{eq:stereovision}
\end{equation}
де $K$ — \emph{інтринсична матриця камери} (два параметри оптичного центру і два параметри масштабування — обидва калібровані), $T$ — перетворення між камерою і глобальними координатами. Дві проекції тієї самої точки можна обчислити прямою тріангуляцією. З $C_L$ як глобальної системи:
\begin{equation}
\begin{pmatrix} u_R\\ v_R\\ 1 \end{pmatrix}
= K T_{LR} \begin{pmatrix} x \\ y \\ z \\ 1 \end{pmatrix}
= K T_{LR} K^{-1} \begin{pmatrix} u_L \\ v_L \\ 1 \end{pmatrix}.
\end{equation}

\emph{Епіполярна лінія}\index{епіполярна лінія}: точка $p_{P_L}$ лежить на промені від центра $C_L$ через точку. Невизначеність — глибина уздовж променя. Проекція цього променя в $P_R$ — епіполярна лінія, вздовж якої треба шукати $p_{P_R}$. Пошук на лінії значно швидший за пошук у всій площині $P_R$.

Виокремлення метрики вимагає унікальної ідентифікації однакових точок. Просте рішення — \emph{структуроване світло}\index{структуроване світло} (рис.~\ref{fig:struclight}). Інфрачервоні сенсори глибини використовують \emph{спекл-патерн} (псевдовипадкові точки з різними відстанями). Ідентифікація однакових точок зводиться до пошуку blob-ів подібного розміру близько один до одного.

\begin{figure}
    \centering
    \includegraphics[width=\textwidth]{figs/structuredlight.png}
    \caption{Зліва направо: два складних об'єкти, патерн кольорових ліній та його деформація на поверхнях, реконструкована 3D-форма. З~\protect\cite{zhang2002rapid}.}
    \label{fig:struclight}
\end{figure}

\section{Комп'ютерний зір і машинне навчання}\label{sec:cvml}
Алгоритми, описані тут, досі формують основу конвеєрів розуміння зображень. З появою згорткових нейромереж (\cref{chap:ann}) базова обробка часто інтегрується в єдиний пайплайн, але розуміння того, що роблять згортка, морфологічні операції, пороги — залишається релевантним.

\section*{Висновки на винос}
\begin{enumerate}
\item На відміну від сенсорів з~\cref{chap:sensors}, наш мозок безпосередньо обробляє 2D-інформацію. Важко ``не думати'' про обробку, яку ми виконуємо автоматично, доповнюючи сигнал знаннями.
\item Описані алгоритми зменшують інформацію до нижньо-вимірного простору, видаляючи шум і допоміжну інформацію.
\item Існує компроміс між трактованістю потоку даних і збереженням реальної інформації. З розвитком обчислень і алгоритмів ML сучасні системи зору об'єднують попередню обробку і розуміння в єдиний конвеєр.
\end{enumerate}

\section*{Вправи}\small
\begin{enumerate}
\item Нижче подано кілька ``ядер'' для згорткового фільтрування зображень.
\begin{equation*}
\begin{array}{|c|c|c|}
\hline
1 & 1 & 1\\
\hline
1 & 2 & 1\\
\hline
1 & 1 & 1\\
\hline
\end{array}
\quad
\begin{array}{|c|c|c|}
\hline
0 & -1 & 0\\
\hline
0 & -1 & 0\\
\hline
0 & -1 & 0\\
\hline
\end{array}
\quad
\begin{array}{|c|c|c|}
\hline
1 & 1 & 1\\
\hline
1 & -4 & 1\\
\hline
1 & 1 & 1\\
\hline
\end{array}
\end{equation*}
\begin{enumerate}
\item Знайдіть ядро, що може розмивати зображення.
\item Які ознаки можна детектувати двома іншими ядрами?
\end{enumerate}
\item Скільки for-циклів потрібно для реалізації 2D-згортки? Поясніть.
\item Використовуйте симулятор робота з камерою у світі з простими об'єктами.
\begin{enumerate}
\item Реалізуйте порогування для виділення об'єкта одного кольору. Чи достатньо простого порога? Чому ні? Чи можна виділити об'єкт нижнім і верхнім порогом?
\item Реалізуйте згладжування: гаусову згортку + морфологічні операції. Експериментуйте з ядрами різної ширини. Переваги/недоліки морфології vs. простий гаусів фільтр?
\item Реалізуйте детектування країв (Sobel). Що ще потрібно для отримання зображення лише з краями?
\end{enumerate}
\item Чи можете придумати алгоритм згладжування, що згладжує малий шум, але зберігає краї? Які фільтри для цього поєднаєте?
\item Знайдіть в інтернеті бібліотеку комп'ютерного зору. Чи реалізує всі алгоритми? Розв'яжіть завдання вище через її функції.
\item Симулюйте дві камери на відомій відстані в одній площині. На простих об'єктах (червоний м'яч) обчисліть їх відстань через стерео-диспаратність.
\end{enumerate} \normalsize
